% page 603 of the facsimile (image p-602 of the scan)
% block 165.1 x 237.7 mm ; left margin 36.9 mm
\pgc{15.240mm}{0.000mm}{
\l{\cel{52}595.}
\l{}
\l{\sou{trilogio}. I. (\sou{epoki antiqua}). Ensemblo de tri tragedii qui}
\l{\cel{3}traktas (en la kazi maxim multa) tri parti intersequanta de}
\l{\cel{3}temo komuna, quin onu reprezentis, ica pos ita, en la kon-}
\l{\cel{3}kursi pri poeziaji qui eventis lor ula festii solena. -}
\l{\cel{3}II. (\sou{nia-epoke}). Asemblajo de tri tragedii qui traktas tri}
\l{\cel{3}parti intersequanta di temo komuna. - DEFIRS.}
\l{}
\l{\sou{trimestro}. Periodo de tri monati, intersequanta. - DEFIRS.}
\l{}
\l{\sou{trimetriko}. (\sou{matem.}) Qua koncernas tri mezuri interdiferanta.}
\l{}
\l{\sou{trioleto}. I. (\sou{poezio}). Kupleto de ok versi en qui la unesma}
\l{\cel{3}repetesas pos la triesma, e la duesma pos la sisesma. -}
\l{\cel{3}II. (\sou{muziko}).Grupo tri-elementa de noti en mezuro du-elementa.}
\l{}
\l{\sou{trinomio}. (\sou{algeb.})\cel{2}Polinomio tri-termina. DEFIRS.}
\l{}
\l{\sou{tripo}. Stomako di la rumineri (pancho, e c.), egardata kom}
\l{\cel{3}nutrivo por la homi.- EFIRS.}
\l{}
\l{\sou{tripliko}. (\sou{yurocienco}). Respondo a dupliko. - F.}
\l{}
\l{\sou{triplikato}. La triesma kopiuro de akto. - DEFIS.}
\l{}
\l{\sou{tripolio.}\cel{1}(\sou{tekn.}) Tero silicoza, flava redatre, uzata por po-}
\l{\cel{3}lisar la speguli, la metali, e c. - DEFIS.}
\l{}
\l{\sou{triptiko}. Pikturo o tabelo sur tri \textquotedbl{}shutri\textquotedbl{} faldebla, de qui}
\l{\cel{3}du (Ia dextra e la sinistra) kovras, pos faldo, la centrala.}
\l{\cel{3}- DEFIS.}
\l{}
\l{\sou{triremo}. (\sou{antique, en Roma}). Navo kun tri rangi de remili. -}
\l{\cel{3}DEFIRS.}
\l{}
\l{\sou{trismo}. (\sou{patol.}) (\sou{che tetananto}). Kontrakturo che la muskuli di}
\l{\cel{3}la maxilo e di la mandibulo. -}
\l{}
\l{\sou{trista}.\cel{2}I. (\sou{pri persono}). Qua sufras psikale. - II. (\sou{pri kozo})}
\l{\cel{3}Qua produktas sufro psikala. - FIS.}
\l{}
\l{\sou{tritono}. I. (\sou{mitol.}) Maro-deajo qua, segun la reprezenti, havas}
\l{\cel{3}vizajo homala e korpo qua finas fishatre. - II. (\sou{zool.})}
\l{\cel{3}Batrako squala, vicina ye salamandro. - L.\cel{1}\sou{triton}. - DEFIRS.}
\l{}
\l{\sou{triturar}. (\sou{trans.})\cel{2}I. Igar pulvero, per frotado, per aplastar}
\l{\cel{3}sen frapar. - II. Igar, de korpo harda, peceti extreme mikra}
\l{\cel{3}per shoki o per la preso da korpo plu harda. - EFIS.}
\l{}
\l{\sou{triumfar}. (\sou{netrans.}) I. (en\cel{1}\sou{la Roma antiqua}). Obtenar la honoro}
\l{\cel{3}qua grantesis a la generalo qua agabis granda vinko, e qua}
\l{\cel{3}konsistis ye enirar pompoze la urbo, sur charo, sequata da}
\l{\cel{3}lua armeo, kun la kaptiti e la raptaji. - II. (\sou{netrans.})}
\l{\cel{3}(ye)\cel{2}Vinkar briloze, komplete. - DEFIRS.}
}
